Now that I have established how software systems can be automatically constructed from meaning using a semantic toolkit, in this section, I will talk about how this drives the future of debugging. 

Earlier in the document, I discussed how shared meaning leads to successful interactions between designs and how that can be used to construct a distributed system. In that example, I highlighted the shared meaning at the boundary of communication between the systems leading to successful interactions. This section expands on that to apply the principle more broadly.

When two humans are unable to communicate effectively to achieve their intentions, the root cause of failure can be identified as an incompatibility in meaning. If an apple means something different to me than it does to another human, there will be a misunderstanding in our communication about apples. If I ask for an apple and they give me a pencil, then the interaction fails. If you were to debug the human interaction, the mismatch in meaning will predict the failure and an agreement on the meaning will resolve the failure.

Now that software systems are defined at the level of meaning using the semantic toolkit and are automatically synthesized as an implementation, the same principles apply. If two software systems share the same meaning, they will not miss understand each other. Another way of saying this is that the system is debugged by constructing it with shared meaning. 

This is not to say that the program can't fail, the shared meaning can be incomplete and a failure will reveal the missing semantics through root cause analysis. However, after learning the new semantics, the new meaning is then shared across to all parts of the system and the semantic toolkit ensures that all the systems are derived from the same meaning. Then through reasoning, the engine can identify all the identify and resolve all the failures. 

This is much more precise than human communication. In human communication, the participants work to ensure that they have the same meaning and human skill plays a part in how effectively the meaning can be remembered and applied. However, in machine communication, they literally share the same meaning. A great analogy would be that it is the same mind realizing its intentions in multiple nodes, it can't misunderstand itself because it is communicating with itself. 

\begin{figure}[htbp]
    \centering
    \includesvg[
        width=0.90\linewidth
    ]{assets/future_debugging.svg}
    \caption{Eliminating bugs within and between systems using semantic toolkit}
    \label{fig:inverted_swe}
\end{figure}

The future of debugging is ensuring that systems understand each other through shared meaning by constructing the meaning of their intentions from the same semantic toolkit and reasoning about their semantics to ensure that all failures in the closed semantic world are eliminated. This process eliminates misunderstanding between and within systems and eliminates failures before they manifest.

Ultimately, this has transformed debugging from a reactive task after a failure happens to a proactive task that eliminates bugs by construction. As a result, systems will not fail in known ways because those failures will be eliminated by construction through shared meaning and reasoning. If they do fail an automatic process will collect the necessary diagnostic information to learn new semantics through root cause analysis and prevent the recurrence of the same failure in the future.

In the next section, in addition to debugging through shared meaning, I will describe how systems of the future will be deterministically intelligent by reasoning about the meaning through their narratives.